\documentclass[11pt,a4paper]{article}
\usepackage[utf8]{inputenc}
\usepackage[T1]{fontenc}
\usepackage{lmodern}
\usepackage[portuguese]{babel}
\usepackage[a4paper,margin=2.4cm]{geometry}
\usepackage{xcolor}
\usepackage{graphicx}
\usepackage{titlesec}
\usepackage{enumitem}
\usepackage{listings}
\usepackage{tcolorbox}
\usepackage{booktabs}
\usepackage{amsmath}
\usepackage{amssymb}
\usepackage{fancyhdr}
\usepackage{hyperref}
\tcbuselibrary{skins,breakable,listings}

\definecolor{BgCream}{HTML}{FAF8F5}
\definecolor{TextEspresso}{HTML}{4A4238}
\definecolor{NudeTaupe}{HTML}{BFAFA6}
\definecolor{SoftBlush}{HTML}{D9C5B2}
\definecolor{SageMuted}{HTML}{7E8F7C}
\definecolor{ClayWarn}{HTML}{B5795C}

\hypersetup{colorlinks=true, linkcolor=TextEspresso, urlcolor=TextEspresso}

\titleformat{\section}{\color{TextEspresso}\Large\bfseries}{\thesection}{0.6em}{}[{\color{NudeTaupe}\titlerule[0.8pt]}]
\titleformat{\subsection}{\color{TextEspresso}\large\bfseries}{\thesubsection}{0.6em}{}
\titlespacing*{\section}{0pt}{1.4\baselineskip}{0.7\baselineskip}
\titlespacing*{\subsection}{0pt}{1.1\baselineskip}{0.5\baselineskip}

\lstset{
  basicstyle=\ttfamily\small,
  breaklines=true,
  columns=fullflexible,
  showstringspaces=false,
  frame=none,
  aboveskip=0pt, belowskip=0pt
}

\pagestyle{fancy}
\fancyhf{}
\renewcommand{\headrulewidth}{0pt}
\fancyfoot[C]{\color{NudeTaupe}\small\thepage}

% ---- Caixa: comandos a executar ----
\newtcblisting{terminal}{
  listing only, breakable,
  colback=BgCream, colframe=NudeTaupe,
  boxrule=0.5pt, arc=2pt, left=6pt, right=6pt, top=5pt, bottom=5pt,
  listing options={basicstyle=\ttfamily\small, breaklines=true, columns=fullflexible}
}

% ---- Caixa: saida esperada ----
\newtcblisting{saida}{
  listing only, breakable,
  colback=white, colframe=SoftBlush,
  boxrule=0.5pt, arc=2pt, left=6pt, right=6pt, top=5pt, bottom=5pt,
  title={\small\bfseries\color{TextEspresso}Saída esperada},
  coltitle=TextEspresso, colbacktitle=SoftBlush!45,
  fonttitle=\small\bfseries,
  listing options={basicstyle=\ttfamily\footnotesize, breaklines=true, columns=fullflexible}
}

% ---- Caixa: ponto de verificacao ----
\newtcolorbox{verificar}{
  breakable, colback=SageMuted!8, colframe=SageMuted,
  boxrule=0.6pt, arc=2pt, left=7pt, right=7pt, top=6pt, bottom=6pt,
  title={\small\bfseries Confirme antes de avançar},
  coltitle=white, colbacktitle=SageMuted, fonttitle=\small\bfseries
}

% ---- Caixa: aviso ----
\newtcolorbox{atencao}{
  breakable, colback=ClayWarn!8, colframe=ClayWarn,
  boxrule=0.6pt, arc=2pt, left=7pt, right=7pt, top=6pt, bottom=6pt,
  title={\small\bfseries Atenção},
  coltitle=white, colbacktitle=ClayWarn, fonttitle=\small\bfseries
}

% ---- Caixa: nota/contexto ----
\newtcolorbox{nota}{
  breakable, colback=NudeTaupe!10, colframe=NudeTaupe,
  boxrule=0.5pt, arc=2pt, left=7pt, right=7pt, top=6pt, bottom=6pt
}

% ---- Cabecalho do manual ----
\newcommand{\manualtitle}[3]{%
  \thispagestyle{empty}
  \begin{center}
  {\color{NudeTaupe}\rule{\textwidth}{0.5pt}}\\[0.5cm]
  {\LARGE\bfseries\color{TextEspresso} Manual Técnico #1}\\[0.35cm]
  {\large\color{TextEspresso} #2}\\[0.3cm]
  {\normalsize\color{NudeTaupe} Infraestruturas de Computação na Nuvem \quad|\quad #3}\\[0.4cm]
  {\color{NudeTaupe}\rule{\textwidth}{0.5pt}}
  \end{center}
  \vspace{0.4cm}
}

% ---- Entrada de troubleshooting ----
\newcommand{\problema}[1]{\vspace{0.35cm}\noindent{\bfseries\color{TextEspresso}#1}\par\nopagebreak\vspace{0.1cm}}

\setlist[itemize]{itemsep=0.22cm, topsep=0.2cm}
\setlist[enumerate]{itemsep=0.22cm, topsep=0.2cm}

\begin{document}
\manualtitle{13}{WBEM com SFCB e Revisão do Semestre}{Aula 13}

\section*{Como usar este manual}

Última aula. Tem três partes: explorar um CIMOM com WBEM, comparar com uma API REST moderna, e rever o semestre.

\begin{atencao}
\textbf{Aviso honesto sobre esta aula.} O SFCB e os seus \emph{providers} são software antigo e pouco mantido, e o empacotamento varia entre versões de distribuição. É possível que algum passo falhe.

Se isso acontecer, \textbf{não perca tempo a insistir}: registe a mensagem de erro exata e peça a alternativa à docente. Analisar e documentar o erro conta como trabalho -- é, aliás, exatamente o que aconteceria num contexto profissional ao lidar com tecnologia legada.
\end{atencao}

\section*{Antes de começar}

\begin{itemize}
  \item[$\square$] Uma VM Linux ligada, com \texttt{sudo}
  \item[$\square$] Os materiais dos guiões anteriores à mão (para a revisão)
  \item[$\square$] Opcional: token de leitura da API do Proxmox, para a Parte 2
\end{itemize}

\section{Parte 1 -- Explorar um CIMOM}

\subsection{Instalar}

\begin{terminal}
$ sudo apt update
$ sudo apt install -y sfcb sblim-cmpi-base cim-schema wbemcli
$ sudo systemctl enable --now sfcb
$ sudo systemctl status sfcb
\end{terminal}

\begin{nota}
Os três pacotes têm papéis distintos: \texttt{sfcb} é o \textbf{CIMOM} (o servidor); \texttt{sblim-cmpi-base} traz os \textbf{providers} que traduzem entre o modelo CIM e o sistema real; \texttt{cim-schema} contém as definições das classes. Sem os providers, o CIMOM corre mas não tem dados para expor.
\end{nota}

Confirme que está a escutar:
\begin{terminal}
$ sudo ss -lntp | grep 5989
\end{terminal}

\subsection{Testar por linha de comandos}

Antes de escrever código, confirme que o CIMOM responde:

\begin{terminal}
$ wbemcli -nl ein 'https://root:<password>@localhost:5989/root/cimv2:CIM_OperatingSystem'
\end{terminal}

\begin{saida}
localhost:5989/root/cimv2:Linux_OperatingSystem.CreationClassName=...,
   CSCreationClassName="Linux_ComputerSystem",CSName="icn-lab",Name="Ubuntu"
\end{saida}

\begin{atencao}
A autenticação é o ponto mais frágil. Se falhar, experimente:
\begin{itemize}
  \item As credenciais do utilizador \texttt{root} local (se tiver password definida)
  \item As suas próprias credenciais de utilizador da VM
  \item Consultar \texttt{/etc/sfcb/sfcb.cfg} para ver o método configurado
\end{itemize}
Se nada funcionar, registe o erro e avance para a Parte 2 -- a docente indicará como obter os dados.
\end{atencao}

\subsection{Instalar o cliente Python}

\begin{terminal}
$ sudo apt install -y python3-venv
$ python3 -m venv ~/wbem-env
$ ~/wbem-env/bin/pip install pywbem
\end{terminal}

\subsection{Enumerar classes}

\begin{terminal}
$ nano explorar_cim.py
\end{terminal}

Copie o script do Guião 13 (substituindo \texttt{PASSWORD} pela sua).

\begin{terminal}
$ ~/wbem-env/bin/python explorar_cim.py | head -30
\end{terminal}

\begin{saida}
=== Classes CIM disponiveis (root/cimv2) ===
CIM_ComputerSystem
CIM_LogicalDisk
CIM_OperatingSystem
CIM_Process
Linux_ComputerSystem
...
\end{saida}

\begin{nota}
Repare nos prefixos: \texttt{CIM\_} são classes do modelo normalizado da DMTF; \texttt{Linux\_} são especializações que herdam das primeiras. É o mecanismo de \textbf{herança} do CIM em funcionamento.
\end{nota}

\subsection{Listar propriedades de uma instância}

Adapte o segundo script do guião e execute-o. Compare com a Aula 12: aqui as propriedades têm \textbf{nomes descritivos e tipos}, em vez de OIDs numéricos.

\subsection{O exercício central: associações}

Esta é a operação que \textbf{não tem equivalente no SNMP}.

\begin{terminal}
$ nano associacoes.py
\end{terminal}

Copie o terceiro script do Guião 13.

\begin{terminal}
$ ~/wbem-env/bin/python associacoes.py
\end{terminal}

\begin{saida}
Sistema: root/cimv2:Linux_ComputerSystem.CreationClassName=...,Name="icn-lab"

=== Associacoes (References) ===
  Linux_CSProcessor
  Linux_RunningOS
  Linux_ComputerSystemProcess

=== Objetos associados (Associators) ===
  Linux_Processor
  Linux_OperatingSystem
  Linux_UnixProcess
\end{saida}

\begin{verificar}
Partindo de \emph{um} objeto (o sistema), descobriu automaticamente os processadores, o sistema operativo e os processos que lhe pertencem -- sem saber previamente que existiam. Registe as associações encontradas e, para cada uma, o que representa no mundo real.
\end{verificar}

\begin{nota}
\textbf{Porque é isto significativo?} Na MIB do SNMP, a relação entre um sistema e os seus discos teria de ser inferida por convenção de índices numéricos. No CIM, a relação é ela própria um objeto que se pode consultar e percorrer. É a diferença entre uma lista e um grafo.
\end{nota}

\subsection{Classes têm métodos}

Execute o quarto script do guião, que mostra as propriedades-chave e os métodos de uma classe.

\begin{verificar}
A classe define \textbf{métodos} -- ações que se podem invocar. No SNMP só existem operações genéricas de leitura e escrita sobre variáveis. É outra diferença estrutural.
\end{verificar}

\subsection{Comparar com o SNMP}

Se ainda tiver o agente da Aula 12 acessível:

\begin{terminal}
$ snmpget -v2c -c icnturma <IP-do-agente> 1.3.6.1.2.1.1.5.0
\end{terminal}

Compare com a forma WBEM de obter a mesma informação. Registe as duas.

\section{Parte 2 -- API REST do Proxmox (opcional)}

\begin{terminal}
$ curl -k -H "Authorization: PVEAPIToken=<user>@pve!<token>=<segredo>" \
    https://<IP-do-servidor>:8006/api2/json/nodes
\end{terminal}

\begin{saida}
{"data":[{"node":"pve-estga","status":"online","maxcpu":32,
  "maxmem":134926512128,"uptime":3721045}]}
\end{saida}

\begin{terminal}
$ curl -k -H "Authorization: PVEAPIToken=..." \
    https://<IP-do-servidor>:8006/api2/json/nodes/<node>/status
\end{terminal}

\subsection{Os três paradigmas lado a lado}

\begin{center}
\begin{tabular}{@{}lll@{}}
\toprule
 & \textbf{Como se identifica} & \textbf{Formato} \\
\midrule
SNMP & OID numérico & Tipos SMI \\
WBEM & Classe + chaves & Objetos CIM \\
REST & URL & JSON \\
\bottomrule
\end{tabular}
\end{center}

\begin{verificar}
Os três resolvem o mesmo problema -- consultar o estado de um sistema remoto -- com 40 anos de distância entre eles. O que mudou foram as ferramentas, não o problema.
\end{verificar}

\section{Parte 3 -- Revisão do semestre}

Percorra os blocos e assinale o que consegue explicar \textbf{sem consultar}. O que ficar por assinalar é o seu plano de estudo para o teste.

\subsection*{C1 -- Datacenters (Aulas 01--02)}
\begin{itemize}
  \item[$\square$] As cinco características essenciais da cloud (NIST)
  \item[$\square$] Distinguir IaaS, PaaS e SaaS, e quem é responsável por quê em cada
  \item[$\square$] Os quatro Tiers e o que distingue o Tier III do Tier II
  \item[$\square$] Calcular disponibilidade em série e em paralelo
  \item[$\square$] Explicar PUE e o que significa um valor próximo de 1
  \item[$\square$] Distinguir tráfego este-oeste de norte-sul
\end{itemize}

\subsection*{C2 -- Recursos computacionais (Aulas 03--04)}
\begin{itemize}
  \item[$\square$] Hypervisor tipo 1 vs tipo 2, com exemplos
  \item[$\square$] \textbf{Porque o contentor partilha o kernel do host e a VM não}
  \item[$\square$] O que fazem os namespaces e o que fazem os cgroups
  \item[$\square$] Explicar thin provisioning e o risco que introduz
  \item[$\square$] Distinguir snapshot, clone completo e template
  \item[$\square$] Contentor de aplicação vs contentor de sistema
\end{itemize}

\subsection*{C3 -- Administração de recursos (Aulas 05--07)}
\begin{itemize}
  \item[$\square$] Os componentes do control plane e o que cada um faz
  \item[$\square$] O ciclo de reconciliação e porque dá self-healing
  \item[$\square$] Para que serve um Service, dado que os Pods são efémeros
  \item[$\square$] \textbf{Porque um nó fica \texttt{NotReady} sem plugin CNI}
  \item[$\square$] Separação entre plano de controlo e plano de dados
  \item[$\square$] Distinguir overlay de underlay
  \item[$\square$] Explicar idempotência com um exemplo concreto
  \item[$\square$] Provisionamento vs gestão de configuração
  \item[$\square$] Para que serve o ficheiro de estado do Terraform
\end{itemize}

\subsection*{C4 -- Fornecimento de serviços (Aulas 08--09)}
\begin{itemize}
  \item[$\square$] Scale-up vs scale-out, e os limites de cada um
  \item[$\square$] Porque um serviço stateless escala mais facilmente
  \item[$\square$] \textbf{Readiness vs liveness probe}
  \item[$\square$] Round-robin vs least connections vs hashing consistente
  \item[$\square$] Camada 4 vs camada 7
  \item[$\square$] Porque o scale-down é mais lento que o scale-up
  \item[$\square$] Redundância ativa-ativa vs ativa-passiva
  \item[$\square$] \textbf{Porque o produtor não é afetado quando o consumidor cai}
  \item[$\square$] Distinguir RTO de RPO
\end{itemize}

\subsection*{C5 -- Armazenamento (Aulas 10--11)}
\begin{itemize}
  \item[$\square$] Block, file e object storage, com um exemplo de cada
  \item[$\square$] RAID 1 vs RAID 5 vs RAID 6
  \item[$\square$] Porque o período de reconstrução é o de maior risco
  \item[$\square$] A pilha PV $\rightarrow$ VG $\rightarrow$ LV
  \item[$\square$] Como funciona um snapshot copy-on-write
  \item[$\square$] \textbf{Porque RAID não substitui backup}
  \item[$\square$] \textbf{Porque dois clientes podem escrever no NFS mas não no mesmo LUN iSCSI}
  \item[$\square$] Enunciar o teorema CAP
\end{itemize}

\subsection*{C6 -- Administração de sistemas (Aulas 12--13)}
\begin{itemize}
  \item[$\square$] Manager, agente, MIB e OID
  \item[$\square$] Counter vs Gauge, e porque o valor absoluto de um contador não chega
  \item[$\square$] \textbf{Porque o SNMPv2c é inseguro e o que o v3 acrescenta}
  \item[$\square$] Polling vs push
  \item[$\square$] O que é fadiga de alertas
  \item[$\square$] \textbf{O que as associações CIM permitem que o SNMP não permite}
  \item[$\square$] Para que serve a gestão fora-de-banda
\end{itemize}

\subsection{Exercício de síntese}

Escolha \emph{um} cenário -- por exemplo, «uma loja online com picos no Natal» -- e escreva um parágrafo por bloco, indicando que tecnologias usaria e porquê.

\begin{nota}
Este exercício não é entregue, mas é a melhor preparação possível para as questões de aplicação do teste, que tipicamente pedem exatamente isto: dado um cenário, propor e justificar uma arquitetura.
\end{nota}

\section{Problemas frequentes}

\problema{\texttt{sfcb.service failed to start}}
Veja a razão concreta:
\begin{terminal}
$ sudo journalctl -u sfcb -n 30 --no-pager
\end{terminal}
Frequentemente falta gerar os certificados:
\begin{terminal}
$ sudo sfcbrepos -f
$ sudo systemctl restart sfcb
\end{terminal}

\problema{\texttt{Authentication failed} no \texttt{wbemcli} ou no \texttt{pywbem}}
O problema mais comum desta aula. Veja o método de autenticação configurado:
\begin{terminal}
$ grep -i auth /etc/sfcb/sfcb.cfg
\end{terminal}
Se estiver \texttt{doBasicAuth: true}, use credenciais válidas do sistema. Se nada funcionar, registe o erro e siga em frente.

\problema{\texttt{EnumerateClassNames} devolve lista vazia}
Os providers não estão instalados ou registados:
\begin{terminal}
$ sudo apt install -y sblim-cmpi-base
$ sudo sfcbrepos -f
$ sudo systemctl restart sfcb
\end{terminal}

\problema{\texttt{Nenhuma instancia de CIM\_ComputerSystem encontrada}}
Nesta instalação a classe pode chamar-se \texttt{Linux\_ComputerSystem}. Liste primeiro as classes disponíveis e ajuste o script ao que existir.

\problema{\texttt{pip install pywbem} falha}
Faltam dependências de compilação:
\begin{terminal}
$ sudo apt install -y build-essential python3-dev libssl-dev
\end{terminal}

\problema{Erro de certificado no \texttt{pywbem}}
Confirme que tem \texttt{no\_verification=True} na criação da \texttt{WBEMConnection}.

\problema{\texttt{curl} à API do Proxmox devolve \texttt{401}}
O formato do cabeçalho é exigente. Tem de ser exatamente:
\begin{terminal}
PVEAPIToken=utilizador@realm!nome-do-token=segredo
\end{terminal}
Confirme que incluiu o \texttt{!} e o realm (\texttt{@pve}).

\section{Referência rápida}

\begin{center}
\begin{tabular}{@{}ll@{}}
\toprule
\textbf{Operação pywbem} & \textbf{Para que serve} \\
\midrule
\texttt{EnumerateClassNames()} & Que classes existem \\
\texttt{EnumerateInstances(cls)} & Instâncias com todas as propriedades \\
\texttt{EnumerateInstanceNames(cls)} & Apenas as referências \\
\texttt{GetClass(cls)} & Definição da classe (propriedades, métodos) \\
\texttt{ReferenceNames(obj)} & Associações que partem do objeto \\
\texttt{AssociatorNames(obj)} & Objetos associados \\
\texttt{InvokeMethod(...)} & Executar um método da classe \\
\midrule
\texttt{sudo sfcbrepos -f} & Reconstruir o repositório de classes \\
\texttt{journalctl -u sfcb -n 30} & Diagnosticar o CIMOM \\
\texttt{wbemcli -nl ein 'https://...'} & Enumerar por linha de comandos \\
\bottomrule
\end{tabular}
\end{center}

\section{Checklist final}

\begin{itemize}
  \item[$\square$] Instalei o SFCB e confirmei que escuta no porto 5989
  \item[$\square$] Enumerei classes e listei propriedades de uma instância
  \item[$\square$] \textbf{Naveguei por associações e registei o que encontrei}
  \item[$\square$] Vi que uma classe CIM define métodos, e não só propriedades
  \item[$\square$] Comparei a forma WBEM com a forma SNMP de obter a mesma informação
  \item[$\square$] (Opcional) Consultei a API REST do Proxmox e comparei os três paradigmas
  \item[$\square$] Percorri a checklist de revisão e identifiquei os pontos a estudar
  \item[$\square$] Fiz o exercício de síntese com um cenário à minha escolha
  \item[$\square$] Confirmei a data de entrega do miniprojeto e a data do teste
\end{itemize}

\end{document}
